% Standalone TeX source. Intentionally not compiled at the user's request.
% Source synthesis and text-level checks: 4 November 2026.
\documentclass[11pt]{article}
\usepackage[margin=1in]{geometry}
\usepackage[T1]{fontenc}
\usepackage[utf8]{inputenc}
\usepackage{lmodern}
\usepackage{amsmath,amssymb,amsthm,mathtools}
\usepackage{booktabs,array}
\usepackage{enumitem}
\usepackage{microtype}
\usepackage[hidelinks]{hyperref}
\usepackage{xurl}
\allowdisplaybreaks
\setlength{\emergencystretch}{2em}
\setlist{itemsep=3pt,topsep=5pt}
\newtheorem{theorem}{Theorem}[section]
\newtheorem{proposition}[theorem]{Proposition}
\newtheorem{lemma}[theorem]{Lemma}
\newtheorem{corollary}[theorem]{Corollary}
\theoremstyle{definition}
\newtheorem{assumption}[theorem]{Assumption}
\newtheorem{definition}[theorem]{Definition}
\theoremstyle{remark}
\newtheorem{remark}[theorem]{Remark}
\newcommand{\R}{\mathbb R}
\newcommand{\E}{\mathbb E}
\newcommand{\Prob}{\mathbb P}
\newcommand{\TV}{\operatorname{TV}}
\newcommand{\Cov}{\operatorname{Cov}}
\newcommand{\Cum}{\operatorname{Cum}}
\newcommand{\diag}{\operatorname{diag}}
\newcommand{\rank}{\operatorname{rank}}
\newcommand{\supp}{\operatorname{supp}}
\newcommand{\dist}{\operatorname{dist}}
\newcommand{\op}{\mathrm{op}}
\newcommand{\tr}{\operatorname{tr}}
\newcommand{\Law}{\mathcal L}
\newcommand{\one}{\mathbf 1}
\newcommand{\trans}{\mathsf T}
\newcommand{\tensor}[2]{#1^{\otimes #2}}
\newcommand{\norm}[1]{\left\lVert#1\right\rVert}
\newcommand{\abs}[1]{\left\lvert#1\right\rvert}
\newcommand{\inner}[2]{\left\langle#1,#2\right\rangle}
\DeclareMathOperator{\Sym}{Sym}

\title{Two Stage Generative Characterization\\[5pt]
\large From conditional quotient discovery to a positive source law\\
and a held out sixth order prediction}
\author{}
\date{4 November 2026}

\begin{document}
\maketitle

\begin{abstract}
Within a declared finite, affine shared-source model, a conditional quotient
and a compatible pair-response chart can become the inputs to a second
experiment. The main result is positive: if an initially unknown centered,
unit-covariance source lies in a known Euclidean ball and an admitted
four-patch law saturates the corresponding fourth-moment bounds in an
orthonormal frame, every compatible source law is the same origin-plus-axes
law in that frame. Its atom weights, positive output generator, and all
higher moments follow. In particular, it predicts a held-out six-patch
score moment. We prove an approximate support and output-law certificate,
including a sharp pointwise rounding constant and an explicit sixth-moment
error bound, and separate population identifiability from calibration and
sampling costs. A complementary fourth-cumulant commutator test excludes
all full-column-rank independent-source explanations of a specified
dimension. It does not certify independence when it passes. The note
states physical access, rank, context, positivity, and temporal assumptions
explicitly, compares the October fixed-law fixture, and credits classical
moment, tensor, and independent component analysis methods. It makes no
literature-priority claim and reports no new physical experiment.
\end{abstract}

\tableofcontents

\section{The result and the boundary of the claim}
\label{sec:scope}

The positive construction has two logically distinct stages:
\begin{enumerate}
\item Conditional single and pair observations discover a task-preserving
quotient and compatible response geometry. They leave a shared frame
ambiguity and do not specify the higher source moments.
\item An actually available simultaneous quartet observation is transformed
using that geometry. Under a separately justified radius bound, fourth
moment saturation determines a positive source distribution and predicts
an unmeasured six-way law.
\end{enumerate}
The first stage supplies a usable statistical representation; it does not
supply the realized latent source or create physical access to a new
experiment. The second stage identifies an entire compatible output-law
family only within the stipulated model and admitted contexts. It does not
identify a unique microscopic mechanism or certify the absence of omitted
physical causes.

The central exact formula is especially simple. If $W\in\R^r$ satisfies
\begin{equation}
\E W=0,\qquad \E WW^\trans=I_r,\qquad \norm{W}_2\le R
\quad\text{almost surely},
\label{eq:normalization-summary}
\end{equation}
and $o_1,\ldots,o_r$ is an orthonormal basis with
$\E(o_j^\trans W)^4=R^2$ for every $j$, then
\begin{equation}
\mu_*=\left(1-\frac r{R^2}\right)\delta_0+
\frac1{2R^2}\sum_{j=1}^r
\bigl(\delta_{Ro_j}+\delta_{-Ro_j}\bigr).
\label{eq:mu-star-summary}
\end{equation}
The coefficient at zero can vanish when $R^2=r$. In the same frame,
$\E(o_j^\trans W)^6=R^4$. If the fourth moment instead has deficit $d_j$,
then
\begin{equation}
0\le R^4-\E(o_j^\trans W)^6\le 2R^2d_j.
\label{eq:sixth-summary}
\end{equation}
These statements require the radius and normalization assumptions; low
order moment matching alone is insufficient.

\subsection{What the predecessor already supplied}

The inspected October incomplete-probe package supplies the following
results, which are inputs or precedents rather than claims of this note:
\begin{itemize}
\item A unique coarsest task-preserving conditional quotient, characterized
by baseline and pair-response block equalities. The strongest displayed
fixture recovers an eight-state alphabet's four-state quotient while
retaining rank-two shared influence on the useful core.
\item A matrix atlas with a single $\mathrm{GL}(r)$ ambiguity on a connected
bipartite observation graph and a common orthogonal ambiguity after a
compatible odd cycle, before joint stochastic positivity is imposed.
\item A positive fixed-source stress family whose two rotated worlds agree
on all pair histories in its specified access model but differ on a
quartet statistic.
\end{itemize}
The separate October two-sided linear evidence result already gives a
fixed-loop criterion, adverse completions, an optional strict inverse
guard, and known-context transport. It is not a theorem of closure for an
arbitrary unknown network, and the present note does not turn it into one.

The difference developed here is the treatment of an \emph{initially
unknown} source law. Rather than assume origin-plus-axes support, the
positive theorem derives it from observed quartet saturation and a known
ball constraint. Section~\ref{sec:october} separates this statement from
what was already built into the earlier fixture.

\section{Model and access assumptions}
\label{sec:model}

Let there be $N$ finite instrumented patches. For patch $i$, let $V_i$ be
its retained quotient output and $c_i$ its pre-update local context,
including any admitted inputs. At one update the same source $W\in\R^r$
is shared across all patches. Write the conditional probability vector as
\begin{equation}
K_i(\cdot\mid c_i,w)=b_i(c_i)+G_i(c_i)w,
\qquad
\one^\trans b_i(c_i)=1,\quad \one^\trans G_i(c_i)=0.
\label{eq:kernel}
\end{equation}
The entries must be nonnegative on the source support. Baselines and
response maps may depend nonlinearly on the local context. Affinity is in
the shared source, not necessarily in local state or control.

\begin{assumption}[Common source and temporal regime]
\label{ass:source}
Fresh patch outputs are conditionally independent given the same current
$W$ and the selected contexts. Source draws are independent and identically
distributed across updates and independent of pre-update contexts. Inputs
precede the current source draw. The unknown law satisfies
\eqref{eq:normalization-summary} for a declared, independently justified
radius $R$. Necessarily $R^2\ge r$ because $\E\norm{W}_2^2=r$.
\end{assumption}
For a single-context population result, the essential source requirement is
the stated common conditional law at that context. Temporal claims and
independent-trial concentration require the additional temporal assumptions.
An unconditional quartet marginal does not identify arbitrary source
memory. Instantaneous output feedback is not implicitly admitted.

\begin{assumption}[Nonempty compatible first stage]
\label{ass:first}
The prior single and pair experiments have certified a nonempty joint
completion family and discovered a conditional quotient. Its relevant
response maps retain rank $r$. The pair data determine a single common
orthonormal source chart, up to a common orthogonal change. A sufficient
setting is a connected nonbipartite full-rank pair observation graph with
jointly compatible factors. Positivity is checked jointly, not independently
at each patch.
\end{assumption}
If the initial graph is bipartite, the pair marginals of a genuinely admitted
full quartet include a same-parity pair and add an odd cycle. That operation
can complete the Euclidean calibration, but its newly used pair information
must be counted as additional evidence. It does not retroactively make the
initial bipartite chart orthogonal.

\begin{assumption}[Actual simultaneous access]
\label{ass:access}
There are four distinct patches whose selected physical contexts are
simultaneously realizable under the same conditional source law. Each of
these four matrices $G_i(c_i)$ has column rank $r$ at that single context.
The operation repertoire permits observing their joint output. Observation
selection must not change the source law. For the held-out sixth-order
experiment, the same requirements hold for six distinct patches, including
conditional independence of all six given the same source occurrence.
\end{assumption}
Full rank after stacking inaccessible or mutually exclusive contexts is not
enough. Separate runs cannot be treated as simultaneous views of one
source draw. Six-view access is an extension of the old four-module
numerical fixture, not something its four-patch result establishes.

\begin{assumption}[Calibrated observations]
\label{ass:calibration}
The conditional output tables and baseline/response representations used
below are observable or have a separately justified calibration and
conditioning argument. If observations are noisy, any sensor inverse and
its conditioning costs are included. Conditioning on hidden old states
requires valid joint old/next-sensor deconvolution, not preparation of a
hidden state or individual decoding of a noisy observation.
\end{assumption}
Exact moment proofs below concern population quantities. Approximate
versions explicitly add calibration and sampling uncertainty.

\section{The first stage produces computable scores}
\label{sec:scores}

\subsection{Conditional quotient discovery}

The predecessor's conditional quotient criterion has a useful algebraic
explanation. For a raw patch response $F_c$, an observed centered pair
table against a full-column-rank reference response $F_{\rm ref}$ has the
form
\begin{equation}
D_c=F_cF_{\rm ref}^\trans
\label{eq:raw-pair}
\end{equation}
after covariance normalization. Multiplication by $F_{\rm ref}^\trans$ is
injective on source row vectors. Therefore equality of summed observed
$D_c$ rows is equivalent to equality of the corresponding summed hidden
response rows. Together with equality of baseline block rows, this tests
pointwise conditional lumpability without choosing a hidden frame.

The prior refinement theorem supplies the unique coarsest task-preserving
quotient under its coverage hypotheses. The order is raw conditional
tables, discovered quotient, aggregated quotient pair tables, and only
then compatible response factors. It is not necessary to supply the
quotient or source coordinates to the discovery step. A retained
conditional quotient need not be autonomous: shared influence can remain
on the useful core.

\subsection{A three-patch chart construction}

At fixed jointly accessible contexts suppress $c_i$ and let
$x_i=e_{V_i}-b_i$, where $e_v$ is the one-hot output vector. The centered
pair laws are
\begin{equation}
D_{ij}=\E[x_ix_j^\trans]=G_iG_j^\trans,\qquad i\ne j.
\label{eq:pair}
\end{equation}
The following complete-graph instance makes the common-chart assumption
explicit.

\begin{proposition}[Pair calibration up to a common orthogonal factor]
\label{prop:chart}
For three distinct full-column-rank patches define
\begin{equation}
H_1=D_{12}(D_{32})^+D_{31}.
\label{eq:gram}
\end{equation}
Then $H_1=G_1G_1^\trans$. Choose any rank-$r$ factor $F_1$ with
$F_1F_1^\trans=H_1$, and set
\begin{equation}
F_i=D_{i1}(F_1^+)^\trans
\label{eq:factor-propagate}
\end{equation}
for each other patch linked to patch $1$. There is one $O\in O(r)$ such
that $F_i=G_iO$ for all these patches.
\end{proposition}
\begin{proof}
For full-column-rank $G_2,G_3$,
$(G_3G_2^\trans)^+=(G_2^\trans)^+G_3^+$.
Substitution in \eqref{eq:gram} cancels the two inner full-rank factors
and yields $G_1G_1^\trans$. Equal full-rank Gram factors differ by a right
orthogonal factor, so $F_1=G_1O$. Since
$(F_1^+)^\trans=G_1(G_1^\trans G_1)^{-1}O$,
\eqref{eq:factor-propagate} gives $F_i=G_iO$.
\end{proof}

A rank-$r$ pair identifies $r$ only within the stipulated
full-column-rank class. More general overcomplete latent representations
are not excluded by this rank statement. The graph version is the
predecessor's matrix-atlas construction, not a new factorization theorem.

Define the recorded-output scores
\begin{equation}
\Phi_i(v)=F_i^+(e_v-b_i).
\label{eq:score}
\end{equation}
They obey
\begin{equation}
\E[\Phi_i(V_i)\mid W]=O^\trans W,
\qquad
\E[\Phi_i(V_i)\Phi_j(V_j)^\trans]=I_r\quad(i\ne j).
\label{eq:score-mean}
\end{equation}
Changing to the learned chart, henceforth write $W$ for $O^\trans W$,
$G_i$ for the learned response factor, and $L_i$ for a left inverse with
$L_iG_i=I_r$. Orthogonal changes preserve the radius bound. A nonorthogonal
chart change generally does not.

\begin{proposition}[Observable moment extraction]
\label{prop:extraction}
For $k$ distinct patches satisfying the common-source conditional
independence and simultaneous-access hypotheses,
\begin{equation}
\E\left[\bigotimes_{i=1}^k L_i(e_{V_i}-b_i)\right]
=\E\tensor{W}{k}.
\label{eq:tensor-extract}
\end{equation}
In particular, a full quartet law determines $T_4=\E\tensor{W}{4}$.
\end{proposition}
\begin{proof}
Conditioning on the common source makes the $k$ score factors independent.
Each has conditional mean $W$. The conditional expectation of their tensor
product is therefore $\tensor{W}{k}$; averaging proves the identity.
\end{proof}

These scores can be negative and large. They are statistical functions of
ordinary outputs, not realizations of $W$, positive preparations, latent
sensors, or intervention controls. Weak responses can make their left
inverses ill-conditioned and moment estimation very expensive.

\section{Fourth moment saturation forces a positive source law}
\label{sec:exact}

\begin{theorem}[Radius saturation]
\label{thm:saturation}
Let $W$ satisfy \eqref{eq:normalization-summary}. Suppose an orthonormal
basis $o_1,\ldots,o_r$ satisfies
\begin{equation}
\E(o_j^\trans W)^4=R^2\quad\text{for every }j.
\label{eq:diagonal-saturation}
\end{equation}
Then its law is exactly \eqref{eq:mu-star-summary}. Conversely that law
satisfies the normalization, radius constraint, and diagonal saturation.
The $r$ diagonal conditions suffice; all mixed fourth moments follow.
\end{theorem}
\begin{proof}
Put $Z_j=o_j^\trans W$. Since $|Z_j|\le R$, the random variable
$Z_j^2(R^2-Z_j^2)$ is nonnegative. Its expectation is
\[
R^2\E Z_j^2-\E Z_j^4=R^2-R^2=0.
\]
It vanishes almost surely, hence $Z_j\in\{0,R,-R\}$ almost surely for
every $j$. The ball bound forbids two nonzero coordinates simultaneously.
Thus $W$ is supported on $0$ and the signed radius-$R$ axes. The variance
identity makes the total mass on the $j$th signed pair $1/R^2$; mean zero
makes its two masses equal. The remaining mass is $1-r/R^2$ at zero.
All these masses are nonnegative. Direct substitution verifies the
converse and all mixed fourth-moment identities.
\end{proof}

\begin{corollary}[Frame and higher moments]
\label{cor:axes}
Under Theorem~\ref{thm:saturation},
\begin{align}
T_4&=R^2\sum_{j=1}^r\tensor{o_j}{4},\label{eq:t4}\\
\E\tensor{W}{2m}&=R^{2m-2}\sum_{j=1}^r\tensor{o_j}{2m}
&& (m\ge1),\label{eq:even}\\
\E\tensor{W}{2m+1}&=0 &&(m\ge0).\label{eq:odd}
\end{align}
The unoriented axis set is identified from $T_4$ up to signed permutation.
\end{corollary}
\begin{proof}
The moment formulas follow by summing over the signed atoms. For unit $u$,
\[
T_4[u,u,u,u]=R^2\sum_j(o_j^\trans u)^4\le R^2,
\]
with equality exactly when one squared coordinate equals one and the
others vanish. Thus the maximizers on the unit sphere are precisely
$\{\pm o_j\}$.
\end{proof}

This is the familiar uniqueness mechanism for an orthogonally decomposable
even-order tensor. Under an unknown-law universe, the response frame and
physical source law are not separately identifiable. Their common-chart
positive mechanism and resulting output laws are the identified objects.

\subsection{The whole-family output generator}

For any admitted simultaneous context tuple $c=(c_1,\ldots,c_k)$, define
\begin{align}
P_*^{(k)}(v\mid c)
={}&\left(1-\frac r{R^2}\right)
\prod_{i=1}^k b_i(v_i\mid c_i)\nonumber\\
&+\frac1{2R^2}\sum_{j=1}^r\sum_{\sigma\in\{-1,1\}}
\prod_{i=1}^k
\left[b_i(v_i\mid c_i)+\sigma R\bigl(G_i(c_i)o_j\bigr)_{v_i}\right].
\label{eq:generator}
\end{align}
All component kernels must be stochastic. In particular, explicitly check
\begin{equation}
b_i(c_i)\pm RG_i(c_i)o_j\ge0
\label{eq:atom-positivity}
\end{equation}
componentwise for every declared context and axis. Exact existence of a
compatible model implies these inequalities for every positive-weight
atom. A reconstructed numerical candidate should still expose them.
The zero kernel is $b_i$, itself an average probability vector. An empty
completion family is rejected, not certified by vacuity.

Under these checks \eqref{eq:generator} is nonnegative, normalized, and
common to every compatible source completion in the learned chart. Its
representation uses at most $2r+1$ mixture components; sampling one
source atom and conditionally independent outputs does not enumerate the
whole global state table. Explicitly listing every $N$-patch output
probability remains exponential in $N$.

\section{A held out sixth order consequence}
\label{sec:sixth}

For six distinct full-rank patches at admitted simultaneous contexts,
Proposition~\ref{prop:extraction} and Corollary~\ref{cor:axes} imply
\begin{equation}
\E\prod_{i=1}^6
\left[o_j^\trans L_i(e_{V_i}-b_i)\right]=R^4,
\qquad j=1,\ldots,r.
\label{eq:six-view}
\end{equation}
In the source-aligned chart $o_j=e_j$. The full sixth tensor, including its
zero mixed entries, is fixed by \eqref{eq:even}. Consequently the entire
six-output law is predicted by \eqref{eq:generator}, not merely its axis
moments.

This is a genuinely held-out implication if one-, two-, and four-patch
laws were used for fitting while the six-patch observations were not.
To execute that validation, six-view access, common-regime context
coverage, rank at each individual context, and six-way conditional
independence must all hold. A table generated from the proposed model is a
regression check, not held-out experimental evidence. No new physical
six-view data have been collected or analyzed for this note.

\section{Approximate support and approximate positive generation}
\label{sec:approx}

Fix an orthonormal candidate frame and use its coordinates $Z$. Retain
exact mean zero, unit covariance, and the exact ball constraint for now.
Suppose
\begin{equation}
\abs{\E Z_j^4-R^2}\le\eta_j,
\qquad
d_j=R^2-\E Z_j^4,
\qquad
\Delta=\sum_{j=1}^r d_j.
\label{eq:deficits}
\end{equation}
The model itself gives $d_j\ge0$, so
$0\le\Delta\le\sum_j\eta_j$. Uniform error $\eta_j\le\eta$ yields
$\Delta\le r\eta$.

Put $s=\sum_j Z_j^2$ and define the nonnegative polynomials
\begin{align}
A(Z)&=s(R^2-s),\nonumber\\
D(Z)&=s^2-\sum_j Z_j^4=\sum_{i\ne j}Z_i^2Z_j^2,\nonumber\\
q(Z)&=A(Z)+D(Z)=R^2s-\sum_j Z_j^4.
\label{eq:defect-polynomials}
\end{align}
Then $\E q=\Delta$, and both $\E A$ and $\E D$ are at most $\Delta$.
Diagonal deficits therefore control their sum directly. Coarser
entrywise fourth-tensor estimates would only give, for example,
$\E A\le r^2\eta$ and $\E D\le r(r-1)\eta$.

\subsection{A sharp rounding inequality}

Let $\mathcal S_R=\{0,\pm Re_1,\ldots,\pm Re_r\}$. Let $U$ be nearest
point rounding of $Z$ to this set, with deterministic tie-breaking. Choose
zero when $\max_j|Z_j|\le R/2$; otherwise choose the signed axis with
largest absolute coordinate.

\begin{lemma}[Sharp distance control]
\label{lem:rounding}
For every $z$ with $\norm{z}_2\le R$,
\begin{equation}
\dist(z,\mathcal S_R)^2\le\frac4{3R^2}q(z).
\label{eq:rounding}
\end{equation}
The constant $4/3$ is sharp for this pointwise inequality.
\end{lemma}
\begin{proof}
Scale to $R=1$. Write $s=\norm{z}_2^2$, $a=\max_j|z_j|$, and note
$\sum_jz_j^4\le a^2s$, hence $q(z)\ge s(1-a^2)$.
If $a\le1/2$, zero is a nearest atom and $q(z)\ge3s/4$.
If $a>1/2$, the nearest squared distance is $1+s-2a$. For
$a^2\le s\le1$, the difference
\[
\frac43s(1-a^2)-(1+s-2a)
\]
is decreasing in $s$, so its minimum is attained at $s=1$ and equals
$\frac23(1-a)(2a-1)\ge0$. Restoring $R$ gives
\eqref{eq:rounding}. At $z=(R/2)e_1$ the inequality is an equality,
so the universal pointwise constant cannot be decreased.
\end{proof}

Consequently, with $e_2=\E\norm{Z-U}_2^2$,
\begin{equation}
e_2\le\frac{4\Delta}{3R^2}\le\frac{4r\eta}{3R^2},
\qquad
\Prob\{\dist(Z,\mathcal S_R)\ge\kappa\}
\le\frac{4\Delta}{3R^2\kappa^2}.
\label{eq:rounding-mean}
\end{equation}
The second assertion holds for every $\kappa>0$. A small positive deficit
does not imply exact finite support.

Two separate geometric bounds are useful if only weaker defect controls
are available. For $0<a\le R^2/2$,
\begin{equation}
\Prob\{a\le s\le R^2-a\}\le\frac{\E A}{a(R^2-a)}.
\label{eq:radial}
\end{equation}
Writing $d_{\rm axis}^2=s-\max_jZ_j^2$, the inequality
$D\ge s\,d_{\rm axis}^2$ gives, for $s_0,\kappa>0$,
\begin{equation}
\Prob\{s\ge s_0,\ d_{\rm axis}\ge\kappa\}
\le\frac{\E D}{s_0\kappa^2}.
\label{eq:angular}
\end{equation}
Both follow from Markov's inequality applied to the indicated nonnegative
polynomials.

\subsection{Weights are controlled as well as support}

Let $\nu=\Law(U)$ and write
\[
t_j=\nu(Re_j)+\nu(-Re_j),\qquad
b_j^{\rm sign}=\nu(Re_j)-\nu(-Re_j).
\]
Use $\mu_*$ for \eqref{eq:mu-star-summary} in this frame and define total
variation as half the $\ell^1$ distance for finite laws.

\begin{proposition}[Approximate atom weights]
\label{prop:weights}
\begin{align}
\sum_j\abs{t_j-R^{-2}}&\le\frac{2\sqrt{e_2}}R,
\label{eq:pair-weights}\\
\sum_j\abs{b_j^{\rm sign}}&\le\frac{\sqrt{re_2}}R,
\label{eq:sign-weights}\\
\TV(\nu,\mu_*)&\le
\min\left\{1,\left(2+\frac{\sqrt r}2\right)
\frac{\sqrt{e_2}}R\right\}.
\label{eq:tv-weights}
\end{align}
\end{proposition}
\begin{proof}
By Cauchy--Schwarz,
\[
\sum_j|U_j^2-Z_j^2|
\le\norm{U-Z}_2\norm{U+Z}_2
\le2R\norm{U-Z}_2.
\]
Taking expectations and using $\E Z_j^2=1$ proves
\eqref{eq:pair-weights}. Mean zero gives
\[
R\sum_j|b_j^{\rm sign}|=\norm{\E U}_1
\le\sqrt r\norm{\E(U-Z)}_2\le\sqrt{re_2},
\]
proving \eqref{eq:sign-weights}. The zero-atom discrepancy is at most the
sum in \eqref{eq:pair-weights}. For each signed pair, its total absolute
probability discrepancy is at most
$|t_j-R^{-2}|+|b_j^{\rm sign}|$. Half the total discrepancy therefore is
at most the first sum plus half the second, yielding
\eqref{eq:tv-weights}.
\end{proof}
Mean zero, independently of the fourth-moment information, fixes the
balance between positive and negative atoms. Fourth moments alone cannot
recover signs.

\begin{theorem}[Approximate positive output law]
\label{thm:tv-output}
Suppose the true affine kernels are stochastic both on the true support
and at every atom of $\mathcal S_R$. For $k$ admitted simultaneous outputs
set $\ell_i=\frac12\norm{G_i}_{2\to1}$. Then
\begin{equation}
\TV(P_{\rm true}^{(k)},P_*^{(k)})
\le\min\left\{1,
\left[\sum_{i=1}^k\ell_i+
\frac{2+\sqrt r/2}{R}\right]\sqrt{e_2}\right\}.
\label{eq:output-tv}
\end{equation}
If estimated stochastic atom kernels have total variation error at most
$\xi_i$ at patch $i$, add $\sum_i\xi_i$ inside the bound, capped at one.
\end{theorem}
\begin{proof}
For source values $z,u$ where the kernels are stochastic,
\[
\TV(K_i(\cdot\mid z),K_i(\cdot\mid u))
=\tfrac12\norm{G_i(z-u)}_1\le\ell_i\norm{z-u}_2.
\]
The total variation distance between product kernels is at most the sum
of the component distances, by coupling or telescoping. Couple $Z$ to
its rounding $U$, average, and use
$\E\norm{Z-U}_2\le\sqrt{e_2}$. Replacing the mixing law of $U$ by
$\mu_*$ costs at most \eqref{eq:tv-weights}, since a stochastic kernel
contracts total variation. A further product-kernel comparison adds the
estimated atom errors.
\end{proof}
Approximate positivity is a separate check. An affine expression that fits
moments is not thereby a probability kernel.

\subsection{Imperfect normalization}

If instead
$\norm{\E Z}_2\le\epsilon_m$ and
$\sum_j|\E Z_j^2-1|\le\epsilon_v$, the weight bound gains
\begin{equation}
\frac{\epsilon_v}{R^2}+
\frac{\sqrt r\,\epsilon_m}{2R}
\label{eq:normalization-error}
\end{equation}
in addition to its rounding term. The proof adds $\epsilon_v/R^2$ to
\eqref{eq:pair-weights} and $\sqrt r\epsilon_m/R$ to
\eqref{eq:sign-weights}. The defect must now be computed as
$\E q=R^2\sum_j\E Z_j^2-\sum_j\E Z_j^4$ with the actual second moments
and their uncertainty. The exact saturation theorem does not apply
unchanged, and a nonorthogonal calibration cannot silently retain the
same ball radius.

\section{Higher moments with explicit error bounds}
\label{sec:higher-approx}

\begin{proposition}[Even and odd axis predictions]
\label{prop:higher}
Under the exact normalization and ball hypotheses used in
Section~\ref{sec:approx}, for $m\ge2$,
\begin{equation}
0\le R^{2m-2}-\E Z_j^{2m}
\le(m-1)R^{2m-4}d_j.
\label{eq:even-error}
\end{equation}
In particular,
\begin{equation}
0\le R^4-\E Z_j^6\le2R^2d_j\le2R^2\eta_j.
\label{eq:sixth-error}
\end{equation}
For $m\ge1$,
\begin{equation}
\abs{\E Z_j^{2m+1}}
\le mR^{2m-2}\sqrt{d_j(R^2-1)}.
\label{eq:odd-error}
\end{equation}
\end{proposition}
\begin{proof}
For $0\le x^2\le R^2$, the difference-of-powers identity gives
\[
0\le R^{2m-2}-x^{2m-2}
\le(m-1)R^{2m-4}(R^2-x^2).
\]
Multiply by $x^2$, take expectations with $x=Z_j$, and use
$\E Z_j^2=1$ to prove \eqref{eq:even-error}. For the odd bound, subtract
$R^{2m}\E Z_j=0$ and factor the difference of powers to obtain
\[
\abs{\E Z_j^{2m+1}}
\le mR^{2m-2}\E\bigl[|Z_j|(R^2-Z_j^2)\bigr].
\]
Cauchy--Schwarz bounds the last expectation by
\[
\sqrt{\E[Z_j^2(R^2-Z_j^2)]\,\E(R^2-Z_j^2)}
=\sqrt{d_j(R^2-1)}.
\]
\end{proof}

Equation~\eqref{eq:sixth-error} concerns the true six-way score moment
under valid exact scores. Validation-table noise and six-view score
calibration must be added before comparing a finite empirical statistic
to this interval. The odd bound expresses approximate symmetry of these
axis moments, not exact symmetry of the source law.

\section{Calibration and finite acquisition for saturation}
\label{sec:calibration}

A good tensor fit alone does not establish a true moment deficit. At the
four selected contexts, let the estimated score be
$\widehat\Phi_i(v)=\widehat L_i(e_v-\widehat b_i)$. Suppose there is one
common true orthogonal alignment $O$ satisfying
\begin{equation}
\norm{\widehat L_iG_i-O}_{\op}\le\kappa_i,
\qquad
\norm{\widehat L_i(b_i-\widehat b_i)}_2\le\beta_i.
\label{eq:calibration-assumptions}
\end{equation}
These are shared-frame calibration guarantees with rank and conditioning
margins, not consequences of a quartet residual. Define
\begin{equation}
a_i=\kappa_iR+\beta_i,
\qquad
B_i^{\rm score}=\max_v\norm{\widehat\Phi_i(v)}_2,
\qquad
C_{\rm cal}=\prod_{i=1}^4(R+a_i)-R^4.
\label{eq:calibration-constants}
\end{equation}
The conditional mean of each estimated score differs from $OW$ by at
most $a_i$. Multilinearity bounds the error in any unit-direction
fourfold product by $C_{\rm cal}$.

If the raw quartet table has $\ell^1$ error at most $\rho$, its transformed
fourth-moment error in any four unit directions is at most
$\rho\prod_i B_i^{\rm score}$. Thus a uniform empirical diagonal fit
error $\eta_{\rm fit}$ yields the valid budget
\begin{equation}
\eta_{\rm eff}=\eta_{\rm fit}
+\rho\prod_{i=1}^4B_i^{\rm score}
+C_{\rm cal}.
\label{eq:eta-effective}
\end{equation}
The norm bounds are uniform over orthonormal rotations, including axes
chosen from the fitted tensor. A probabilistic claim still requires a
simultaneous confidence event for the raw table, baselines, and shared
frame. Conditional tables obtained by dividing joint frequencies require
context-mass floors, with their uncertainty and any deconvolution errors
included in these budgets.

\subsection{Independent validation after freezing the fitted axes}

One implementable confidence check uses sample splitting. Fit the shared
chart and candidate axes on training data and freeze them. On $n$
independent validation quartet windows at the required contexts, each
axis score product lies in $[-M,M]$, where
$M=\prod_i B_i^{\rm score}$. Hoeffding's inequality and a union bound over
$r$ axes give simultaneous sampling error at most $t$ with probability
at least $1-\alpha$ when
\begin{equation}
n\ge\frac{2M^2}{t^2}\log\frac{2r}{\alpha}.
\label{eq:saturation-hoeffding}
\end{equation}
A valid upper bound on the true deficit along frozen axis $j$ is
\begin{equation}
\overline d_j=
\max\{0,R^2-\widehat T_j+t+C_{\rm cal}\}
\label{eq:deficit-confidence}
\end{equation}
on the joint calibration and validation event, where $\widehat T_j$ is
the empirical diagonal product. Add the calibration failure probability
to $\alpha$. If a confidence interval contradicts the model-imposed upper
bound $R^2$, report inconsistency rather than hiding it by clipping.

This calculation does not establish a finite-sample pair-factorization
algorithm or infer a physical support radius from data. Dependent
trajectory windows need a valid replacement concentration bound. No
finite $n$ proves an exact zero deficit. Small responses can make $M$ so
large that a formally identified experiment is impractical.

\section{Why the radius and saturation conditions are essential}
\label{sec:counterexamples}

\subsection{An unsaturated scalar example}

Consider the symmetric scalar laws
\begin{align*}
\mu_A&=\tfrac12\delta_0+
\tfrac14(\delta_{\sqrt2}+\delta_{-\sqrt2}),\\
\mu_B&=\tfrac13\delta_0+
\tfrac14(\delta_1+\delta_{-1})+
\tfrac1{12}(\delta_{\sqrt3}+\delta_{-\sqrt3}).
\end{align*}
Both lie in $[-\sqrt3,\sqrt3]$, are centered, and have variance one and
fourth moment two. Their sixth moments are four and five respectively.
Here $R^2=3$, so the common fourth moment is unsaturated. In any affine,
conditionally independent output model with the same valid channels,
all laws through four distinct patches agree, but a six-patch score can
differ. Finite low-order moments generally do not determine the source.

\subsection{A growing-rank counterexample when the radius is loosened}

For any $r\ge1$, the target law has radius squared $R^2=4r/3$, mass
$1/4$ at zero, and mass $3/(8r)$ at each of the $2r$ signed radius-$R$
axes. Compare the alternative law
\begin{equation}
\widetilde\mu=\frac16\delta_0+
\frac1{3r}\sum_{j=1}^r(\delta_{\sqrt r e_j}+\delta_{-\sqrt r e_j})+
\frac1{12r}\sum_{j=1}^r
(\delta_{\sqrt{2r}e_j}+\delta_{-\sqrt{2r}e_j}).
\label{eq:radius-alternative}
\end{equation}
Both are centered with covariance $I_r$, every fourth diagonal $4r/3$,
and all mixed fourth moments zero. All moments through order four agree.
Their sixth diagonals are respectively $16r^2/9$ and $2r^2$.
The alternative needs radius $\sqrt{2r}$, which violates the target's
radius $\sqrt{4r/3}$. For $r=2$, the sixth diagonals are $64/9$ and $8$,
and the alternative radius is $2$.

These are positive finite-output alternatives, not just abstract moment
lists. Use $2r$ output symbols, baseline $1/(2r)$, and response rows
$\pm\epsilon e_j^\trans$ with $\epsilon=1/(8r^2)$. The response has rank
$r$ and all output probabilities are positive on the larger ball, since
$\epsilon\sqrt{2r}<1/(2r)$. Conditional independence then makes the full
output laws agree for up to four patches, while a six-view score
separates them. Thus the independently justified smaller radius performs
real identifying work.

\section{Comparison with the October fixture}
\label{sec:october}

The October stress family \emph{declared} its full source law:
\begin{equation}
Z\in\{0,\pm e_1,\pm e_2\},\qquad
\Prob(Z=0)=1-\tau,\qquad
\Prob(Z=\pm e_j)=\frac\tau4.
\label{eq:old-law}
\end{equation}
At $\tau=3/4$, its covariance is $cI_2$ with $c=3/8$. Standardizing
$W=Z/\sqrt c$ gives
\begin{equation}
r=2,\qquad R^2=\frac83,\qquad
\E W_j^4=\frac83,\qquad \E W_j^6=\frac{64}9.
\label{eq:old-values}
\end{equation}
The old full-column-rank four-row contrast matrix is
\begin{equation}
U=\begin{pmatrix}1&0\\-1&0\\0&1\\0&-1\end{pmatrix}.
\label{eq:old-U}
\end{equation}
Its response has the form $s_i(c)UA_i$ with $s_i(c)>0$ and $A_i$
invertible, so rank two holds at each displayed individual context,
not merely after stacking contexts.

The current premise weakens \eqref{eq:old-law} to fresh IID draws, mean
zero, covariance $cI_2$, and $\norm{Z}\le1$, while allowing an initially
unknown distribution. Rotated-axis and continuous laws, among many
others, satisfy those preliminary constraints. Quartet saturation then
forces \eqref{eq:old-law} in the learned response chart. It is incorrect
to describe this as discovering an unknown support in the original
fixed-law fixture: the original had already assumed it. The conditional
positive result concerns the broader bounded-law model.

The bound remains strong and external. If only a larger slack radius is
known, the saturation certificate can fail even when the true law happens
to be the old five-point law. The radius counterexample of
Section~\ref{sec:counterexamples} shows why that distinction is necessary.

The prior quartet test distinguished two named rotated worlds. The
positive theorem instead determines all compatible source distributions
under its strengthened moment evidence and radius hypothesis. Its
sixth-order prediction extends the numerical fixture to six distinct
patches; it does not report a physical six-patch experiment.

\section{A complementary fourth cumulant exclusion test}
\label{sec:commutator}

The positive saturation result uses a bounded-source promise. A separate
construction needs only mean zero, unit covariance, and finite fourth
moments, together with the same affine, full-rank, simultaneous four-view
model. Source coordinates are \emph{not} assumed independent.

Let $z_i=L_i(e_{V_i}-b_i)$ be exact scores in the learned chart. For
$g\in\R^r$ collect a quartet and compute
\begin{equation}
W_g=z_1z_2^\trans(g^\trans z_3)(g^\trans z_4),
\qquad
M_g=\E W_g-\norm g_2^2I_r-2gg^\trans.
\label{eq:commutator-stat}
\end{equation}
The subtractions are the three cross-pair products, using distinct
patches rather than repeated same-patch measurements. The population
matrix is symmetric, with
\begin{equation}
(M_g)_{ab}=\Cum(W_a,W_b,g^\trans W,g^\trans W).
\label{eq:cumulant-contraction}
\end{equation}
It can be symmetrized empirically. Two $r\times r$ summaries suffice for
a two-direction test; the full categorical quartet table or an $r^4$
tensor need not be stored.

Under a common orthogonal chart change and transported directions,
$M_g$ changes by orthogonal conjugation. Consequently vanishing and
operator norm of $[M_g,M_h]=M_gM_h-M_hM_g$ are invariant. Reusing the
same numeric vectors in a different chart may select different intrinsic
directions, but any nonzero result remains a valid exclusion certificate.

\begin{theorem}[Exclusion of full-rank independent-source realizations]
\label{thm:exclusion}
If a full-column-rank, $r$-dimensional independent-source realization of
the observed singles, pairs, and quartet statistics exists in the declared
affine common-source class, then $[M_g,M_h]=0$ for every $g,h$.
One nonzero measured population commutator therefore excludes every
such realization, not merely one chosen orientation.
\end{theorem}
\begin{proof}
Standardize each independent source coordinate to variance one, absorbing
its scale into its response column. Its fourth-cumulant tensor is diagonal
in the independent basis. In the learned chart it has the form
\begin{equation}
\mathcal C=\sum_{j=1}^r\kappa_j\tensor{q_j}{4}
\label{eq:odeco-cumulant}
\end{equation}
for an orthonormal basis $q_j$ and fourth cumulants $\kappa_j$. Thus
\begin{equation}
M_g=\sum_{j=1}^r\kappa_j(q_j^\trans g)^2q_jq_j^\trans.
\label{eq:diagonal-contraction}
\end{equation}
All such matrices are diagonal in the same orthonormal basis and commute.
Every compatible full-rank pair factorization differs by one common
orthogonal chart change, so another latent frame cannot evade the test.
\end{proof}
A larger overcomplete independent generator without these rank promises
is outside this exclusion statement. The rejection theorem does not
require an independence promise; it tests that stronger proposed
explanation within the common affine-source universe.

\subsection{The exact algebraic boundary is only fourth order}

\begin{proposition}[Commuting contractions and orthogonal decomposition]
\label{prop:odeco}
For a symmetric fourth-order tensor $\mathcal C$, all matrices
$\mathcal C(\cdot,\cdot,g,g)$ commute if and only if it has the form
\eqref{eq:odeco-cumulant}, allowing zero coefficients.
\end{proposition}
\begin{proof}
One direction is \eqref{eq:diagonal-contraction}. Conversely, the
finite-dimensional span of commuting symmetric contractions admits a
common orthonormal eigenbasis. In that basis,
$\mathcal C_{ijgg}=0$ for every $g$ when $i\ne j$. Polarization implies
$\mathcal C_{ijkl}=0$ for every $k,l$ whenever $i\ne j$. Full tensor
symmetry can move any two unequal indices into the first two positions,
so every entry except $\mathcal C_{iiii}$ vanishes.
\end{proof}

An exact finite test uses $g=e_i$ and $g=e_i+e_j$, whose matrices $gg^\trans$
span $\Sym^2(\R^r)$, and checks all pairs of their contractions.
Alternatively, if the tensor is not orthogonally decomposable, a pair
$(g,h)$ drawn from an absolutely continuous distribution yields a nonzero
commutator almost surely: at least one entry is a nonzero polynomial in
their coordinates. This population statement does not imply a reliable
finite-noise acceptance test.

Passing the test does not establish statistical independence. For $r\ge5$,
the uniform law on sign vectors with $\prod_jW_j=1$ is dependent but has
the same moments through order four as independent Rademacher coordinates.
It has $2^{r-1}$ support points. Any nondegenerate independent law on $r$
coordinates has at least $2^r$ support points, so an invertible rotation
cannot turn this law into an independent $r$-source law. A fourth-order
procedure cannot distinguish these two moment tensors.

\section{A positive categorical family in every rank}
\label{sec:family}

For $r\ge2$ set $q=r+1$ and choose a $q\times r$ matrix $U$ with
$U^\trans U=I_r$ and $U^\trans\one=0$. Choose arbitrary unknown orthogonal
matrices $R_1,\ldots,R_4$ and define
\begin{equation}
G_i=\lambda UR_i,\qquad
K_i(a\mid w)=\frac1q+\lambda(UR_iw)_a,
\qquad
0<\lambda\le\frac1{2(r+1)\sqrt r}.
\label{eq:positive-family}
\end{equation}
For source values of norm $\sqrt r$, each entry is at least $1/(2q)$,
since each row of $U$ has norm at most one. The entries sum to one and
each response has rank $r$.

Compare these sources using the same four channels:
\begin{enumerate}
\item $I_r$: the uniform law on $\{-1,1\}^r$.
\item $C_r$: with probability $r/(r+2)$ use $I_r$, and with probability
$2/(r+2)$ use the uniform law on $\{\pm\sqrt r e_j:1\le j\le r\}$.
\end{enumerate}
Both are centered, have covariance $I_r$, and have zero third moments by
global sign symmetry. A product of at most three affine channel
probabilities depends only on source moments through degree three.
Therefore every single, pair, and triple output law agrees exactly.

For $C_r$,
\[
\E W_i^4=\frac{3r}{r+2},\qquad
\E W_i^2W_j^2=\frac r{r+2}\quad(i\ne j),
\]
and every fourth moment with an odd coordinate multiplicity is zero.
With $\beta=-2/(r+2)$, its fourth-cumulant tensor and contractions are
\begin{align}
\mathcal C_{ijkl}
&=\beta(\delta_{ij}\delta_{kl}+\delta_{ik}\delta_{jl}
+\delta_{il}\delta_{jk}),\nonumber\\
M_g&=\beta(\norm g_2^2I_r+2gg^\trans).
\label{eq:coupled-cumulant}
\end{align}
The tensor is rotationally invariant although the full source law need
not be. In any learned chart, take $g=e_1$ and
$h=(e_1+e_2)/\sqrt2$. Direct multiplication gives
\begin{equation}
[M_g,M_h]=2\beta^2(e_1e_2^\trans-e_2e_1^\trans),
\qquad
\norm{[M_g,M_h]}_{\op}=\frac8{(r+2)^2}.
\label{eq:gap}
\end{equation}
For $I_r$ the cumulant tensor is $-2\sum_j\tensor{q_j}{4}$ in the learned
chart and all contractions commute. The quartet experiment thus excludes
every full-rank independent-source explanation of $C_r$, not only the
displayed $I_r$ comparison. The dependence involves every pair of source
directions rather than appending independent coordinates to a fixed
low-dimensional example.

The mixture has $2^r+2r$ support points but an $O(r)$ sampler: choose the
branch, then independent signs or one coordinate and a sign. Summary and
experimental dimensions grow polynomially. This is not a claim of an
economical sample requirement.

\section{Noise bounds for the exclusion test}
\label{sec:commutator-noise}

\begin{proposition}[Sound one-sided rejection]
\label{prop:noise}
Suppose $\norm{\widehat M-M_g}_{\op}\le\epsilon$ and
$\norm{\widehat N-M_h}_{\op}\le\epsilon$. Then
\begin{equation}
\norm{[\widehat M,\widehat N]-[M_g,M_h]}_{\op}
\le2\epsilon(\norm{\widehat M}_{\op}+\norm{\widehat N}_{\op})
+2\epsilon^2.
\label{eq:commutator-bound}
\end{equation}
Consequently the strict computable inequality
\begin{equation}
\norm{[\widehat M,\widehat N]}_{\op}
>2\epsilon(\norm{\widehat M}_{\op}+\norm{\widehat N}_{\op})
+2\epsilon^2
\label{eq:rejection}
\end{equation}
certifies a nonzero true commutator on that error event.
\end{proposition}
\begin{proof}
Write $M_g=\widehat M-E$ and $M_h=\widehat N-F$. Expand their commutator
and use $\norm{[A,B]}_{\op}\le2\norm A_{\op}\norm B_{\op}$ on the two
linear error terms and the quadratic term.
\end{proof}
Failure of \eqref{eq:rejection} is inconclusive, not a certificate of
independence.

\subsection{Exact score calibration}

To avoid confusing a source radius with a score bound, write $H$ for a
uniform bound $\norm{z_i}\le H$. For unit $g,h$, every entry of a
symmetrized $W_g$ or $W_h$ lies in $[-H^4,H^4]$. With $n$ independent
quartet trials, Hoeffding, a union bound over $2r^2$ entries, and
$\norm E_{\op}\le r\max_{ab}|E_{ab}|$ give both matrix errors at most
$\epsilon$ with probability at least $1-\alpha$ if
\begin{equation}
n\ge\frac{2H^8r^2}{\epsilon^2}\log\frac{4r^2}{\alpha}.
\label{eq:commutator-hoeffding}
\end{equation}
The same quartet sample estimates both contractions. No separate
physical operation is required for each tensor entry.

For \eqref{eq:positive-family}, $H\le1/\lambda$. In $C_r$, the true
contraction norms are $3|\beta|$. Set
$\epsilon=|\beta|/16=1/[8(r+2)]$. A direct expansion around the true
matrices bounds the commutator error by
$12\epsilon|\beta|+2\epsilon^2=(97/128)\beta^2$. Hence the empirical
commutator norm is at least $(159/128)\beta^2$, while the threshold in
\eqref{eq:rejection} is at most $(99/128)\beta^2$. Rejection is strict.
At $\lambda=1/[2(r+1)\sqrt r]$, a sufficient sample size is
\begin{equation}
n\ge32768(r+1)^8r^6(r+2)^2\log\frac{4r^2}{\alpha}.
\label{eq:commutator-sample-large}
\end{equation}
This $O(r^{16}\log(r/\alpha))$ bound is deliberately conservative, not
optimal and not a practical sample recommendation. It exposes the price
of weak categorical readouts.

\subsection{Estimated pair calibration}

The preceding sample count covers quartet acquisition only. Pair samples,
singular-value conditioning, and calibration confidence are additional.
After one common orthogonal alignment, suppose
\[
\norm{\widehat z_i(a)-z_i(a)}_2\le\eta\quad\text{for every outcome }a,
\qquad \norm{z_i(a)}_2\le H.
\]
Using independent quartet data and the same ideal-chart Wick subtraction,
telescoping the four factors bounds deterministic contraction error by
\begin{equation}
\epsilon_{\rm cal}=(H+\eta)^4-H^4
\le4\eta(H+\eta)^3.
\label{eq:commutator-cal}
\end{equation}
Use $H+\eta$ in the sampling bound and
$\epsilon=\epsilon_{\rm stat}+\epsilon_{\rm cal}$ in
\eqref{eq:rejection}. Add calibration failure probabilities to the
statistical failure probability. Baseline error belongs in $\eta$.
An honest $\eta$ requires rank and conditioning margins and a valid
calibration confidence argument. It does not follow from finite data or
a small quartet residual alone. Hidden source memory requires separate
treatment.

\section{Conditional axis recovery and the ICA lineage}
\label{sec:ica}

Under the additional promise that coordinates are independent and at most
one fourth cumulant is zero, a generic contraction identifies their axes.
Indeed \eqref{eq:diagonal-contraction} has eigenvalues
$\kappa_j(q_j^\trans g)^2$. For any pair of coordinates, equality of these
eigenvalues is a nonzero polynomial equation unless both cumulants are
zero. An absolutely continuous random $g$ therefore produces a simple
spectrum almost surely. Equal nonzero kurtoses do not obstruct this
argument. Its eigenvectors recover axes up to signed permutation.

The recovered $G_iq_j$ are response directions, not observed source
values or controls. If the simple-spectrum gap is $\gamma$ and the
symmetric contraction error is $\epsilon<\gamma/4$, a standard
symmetric eigenspace perturbation bound gives, after matching,
$\sin\angle(\widehat q_j,q_j)\le2\epsilon/\gamma$.
Calibration error must be included in $\epsilon$.

This is classical fourth-cumulant independent component analysis after
pair-chart construction. An independence promise cannot be inferred from
the first-stage data or from commuting fourth-order contractions. The
classical ICA and JADE lineage is essential context, not an incidental
comparison; see Miettinen et al.~\cite{miettinen}. Algebraic
characterizations of orthogonal tensor decompositions are developed by
Boralevi et al.~\cite{boralevi}. Multiview ICA with shared and individual
sources has its own substantial literature, including
Pandeva and Forr\'e~\cite{pandeva}; their noisy linear observation model
is not identical to this categorical conditional-mean model.

\subsection{The old five-point law already fails the independence test}

For the standardized October law \eqref{eq:old-law},
\[
\mathcal C_{1111}=\mathcal C_{2222}=2/\tau-3,
\qquad
\mathcal C_{1122}=\mathcal C_{1212}=\mathcal C_{1221}=-1,
\]
with odd-coordinate fourth entries zero. At $\tau=3/4$, in the displayed
source-aligned frame,
\begin{align}
K_{11}&=\mathcal C(\cdot,\cdot,e_1,e_1)
=\begin{pmatrix}-1/3&0\\0&-1\end{pmatrix},\nonumber\\
K_{12}&=\mathcal C(\cdot,\cdot,e_1,e_2)
=\begin{pmatrix}0&-1\\-1&0\end{pmatrix},\nonumber\\
[K_{11},K_{12}]&=\begin{pmatrix}0&-2/3\\2/3&0\end{pmatrix},
\qquad \norm{[K_{11},K_{12}]}_{\op}=\frac23.
\label{eq:old-commutator}
\end{align}
The mixed contraction is observable by polarization:
$K_{12}=(M_{e_1+e_2}-M_{e_1}-M_{e_2})/2$.
Equivalently $M_{(e_1+e_2)/\sqrt2}=-(2/3)I+K_{12}$, so it has the
same commutator norm with $K_{11}$.

Pair observations alone do not supply this numeric source-aligned frame.
In an arbitrary learned chart use the spanning-set test or generic
random directions. Non-orthogonal-decomposability is frame independent,
though the displayed entries are not. Both of the old rotated worlds
fail the independence criterion; it does not replace the old quartet
statistic for distinguishing those orientations. Since the old fixture
specified its whole law, its independence failure was already a
consequence of its structural input. The substantive unknown-source
exclusion experiment concerns the broader affine-source class.

\section{Verification and interpretation}
\label{sec:verification}

\subsection{Exact rational support checks}

The companion support script uses Python rational arithmetic with exact
equality assertions. Its recorded result is \texttt{PASS}. It checks:
\begin{itemize}
\item The old $\tau=3/4$ values $c=3/8$, $R^2=8/3$, fourth diagonal
$8/3$, and sixth diagonal $64/9$.
\item Rank two at one physical context, a positive $4^4=256$-cell quartet
table, and all $2^4=16$ extracted fourth tensor entries.
\item A model-generated $4^6=4096$-cell six-view table, its axis sixth
moments, and a mixed sixth moment.
\item Equal single, pair, and triple laws for a rotated source, alongside
distinct quartet laws.
\item An approximate rotated mixture, the rounding inequality on 1257
rational grid points, and its equality case.
\item The scalar unsaturated example and the radius counterexamples for
$r=1,2,3,8,16,64$.
\end{itemize}
The recorded approximate-mixture defect $108/3125$ and mean squared
rounding distance $3/100$ use the script's \emph{raw unit-ball} coordinates
with covariance $3I/8$, not the standardized coordinates of
Section~\ref{sec:approx}. They must not be substituted into standardized
formulas without the corresponding scaling.

These are exact arithmetic regressions for finite constructions, not a
formal verification of the universal theorems, a sampled-data calibration
study, or physical held-out validation. The six-view table is generated
from the model.

\subsection{Floating point commutator checks}

The separate executable constructs actual categorical joint laws for
$r=2,3,4,6,8$ with unrelated random patch rotations, reconstructs the
chart from pair laws, and contracts each actual $q^4$ quartet law with
the learned scores. Its recorded gaps are summarized below.

\begin{center}
\begin{tabular}{@{}rrr@{}}
\toprule
$r$ & Coupled commutator norm & Theoretical value\\
\midrule
2 & $0.500000000000001$ & $1/2$\\
3 & $0.320000000000023$ & $8/25$\\
4 & $0.222222222222177$ & $2/9$\\
6 & $0.125000000000091$ & $1/8$\\
8 & $0.080000000000412$ & $2/25$\\
\bottomrule
\end{tabular}
\end{center}
The largest recorded pair or triple discrepancy is below
$3\times10^{-17}$; independent-world commutator norms are below
$7\times10^{-13}$. These are floating-point consistency checks, not
formal proofs, human peer review, or finite-sample validation.

\subsection{What is and is not completed}

The mathematical work provides a conditional exact theorem, approximate
bounds, a positive generative representation, an explicit held-out
prediction, and reproducible model checks. The operational validation
still requires justified source-radius and temporal hypotheses, admissible
four- and six-view access, context coverage, calibrated response maps,
joint positivity, and collected data with valid uncertainty bounds.

No theorem here discovers an arbitrary unknown patch cover, unrestricted
nonlinear hidden mechanism, or general cross-module causal architecture.
No moment calculation supplies physical ownership or control of a latent
coordinate. A compatible statistical generator is not a unique physical
ontology. The construction is a bounded example of a representation
learned in one stage enabling a new, explicitly admitted experiment in
the next.

\section{Prior work and source records}
\label{sec:sources}

The orthogonal tensor moment methodology has a classical baseline in
Anandkumar et al.~\cite{anandkumar}. The support proof is elementary
moment-problem reasoning: the support of a representing probability
measure lies in the zero set of a nonnegative polynomial whose integral
vanishes. Related moment-matrix kernel/support statements appear in
Curto and Fialkow~\cite{curto}. Neither the saturation proof nor its use
here is presented as a new general moment theorem. The commutator
construction uses the established fourth-cumulant ICA and orthogonal
tensor lineage discussed in Section~\ref{sec:ica}.

Relative to the inspected project package, this note makes explicit the
bridge from a discovered conditional quotient and compatible pair frames
to an admitted quartet observation, a bounded-law support and weight
certificate, a positive output generator, and a held-out sixth-order
prediction with quantitative error bounds. This describes a project-level
continuation, not a literature-priority claim. No exhaustive novelty
assessment is asserted.

\subsection{Project source records}

The main mathematical sources synthesized here are:
\begin{itemize}
\item \emph{Two-stage continuation: discovered components to a forced
shared-source law}, 4 November 2026,
\path{two_stage_setting_audit/Quartic_Support_Continuation.md}.
\item \emph{From a pair-observation chart to an executable independence
test}, 4 November 2026,
\path{generativity_independence_study_20261104.md}.
\end{itemize}
The inspected predecessor records include
\path{incomplete_probe/main/Vector_Directional_Theorem.md},
\path{incomplete_probe/gluing/Matrix_Atlas_Classification.md}, and
\path{incomplete_probe/stress/Matrix_Gauge_Family.md}.
The latter's conditional discovery fixture and fixed-source quartet test
are the relevant October comparisons.

The accompanying check records are
\path{two_stage_setting_audit/check_quartic_support.py},
\path{two_stage_setting_audit/check_quartic_support.json},
\path{generativity_commutator/check.py}, and
\path{generativity_commutator/results.json}.
They are provenance references rather than TeX build dependencies. This
source file is standalone and contains no external figures, bibliography
files, or included TeX fragments. It has intentionally not been compiled.

\begin{thebibliography}{9}

\bibitem{anandkumar}
A.~Anandkumar, R.~Ge, D.~Hsu, S.~M.~Kakade, and M.~Telgarsky.
\newblock Tensor decompositions for learning latent variable models.
\newblock \emph{Journal of Machine Learning Research}, 15:2773--2832, 2014.
\newblock \url{https://www.jmlr.org/papers/v15/anandkumar14b.html}.

\bibitem{curto}
R.~E.~Curto and L.~A.~Fialkow.
\newblock Solution of the singular quartic moment problem.
\newblock See the kernel-polynomial support relation at equation (1.2).
\newblock \url{https://www.cs.newpaltz.edu/~fialkowl/files/tcmp6.pdf}.

\bibitem{miettinen}
J.~Miettinen, S.~Taskinen, K.~Nordhausen, and H.~Oja.
\newblock Fourth moments and independent component analysis.
\newblock \emph{Statistical Science}, 30(3):372--390, 2015.
\newblock \url{https://arxiv.org/abs/1406.4765}.

\bibitem{boralevi}
A.~Boralevi, J.~Draisma, E.~Horobe\c{t}, and E.~Robeva.
\newblock Orthogonal and unitary tensor decomposition from an algebraic
perspective.
\newblock \emph{Israel Journal of Mathematics}, 222(1), 2017;
preprint 2015.
\newblock \url{https://arxiv.org/abs/1512.08031}.

\bibitem{pandeva}
T.~Pandeva and P.~Forr\'e.
\newblock Multi-view independent component analysis with shared and
individual sources.
\newblock \emph{Proceedings of the Thirty-Ninth Conference on Uncertainty
in Artificial Intelligence}, PMLR 216:1639--1650, 2023.
\newblock \url{https://proceedings.mlr.press/v216/pandeva23a.html}.

\end{thebibliography}
\end{document}
